\documentclass[12pt]{letter}
\usepackage[utf8]{inputenc}
\usepackage{geometry}
\geometry{margin=1in}

\signature{Min Woo Park\\Independent Researcher\\nerve011235@gmail.com}
\address{Min Woo Park\\Independent Researcher\\Republic of Korea}
\date{March 30, 2026}

\begin{document}

\begin{letter}{Editor-in-Chief\\The American Mathematical Monthly\\Mathematical Association of America}

\opening{Dear Editor,}

I am pleased to submit the manuscript entitled ``Perfect Number 6 Cross-Domain Unification: Arithmetic Functions as Architecture Predictors'' for consideration in \textit{The American Mathematical Monthly}.

This paper demonstrates an unexpected connection between the arithmetic functions of the smallest perfect number, 6, and the empirically optimal parameters of neural network architectures. We show that $\tau(6)=4$, $\sigma(6)=12$, $\varphi(6)=2$, and the divisor set $\{1,2,3,6\}$ each predict independently verified architectural constants---attention head counts, layer depths, initial cell counts, and developmental stages---across models ranging from 4M to 700M parameters.

I believe this work is well suited to \textit{The American Mathematical Monthly} for the following reasons:

\begin{enumerate}
    \item It presents a novel, accessible application of classical number theory (perfect numbers, divisor functions, Euler's totient) to a modern domain.
    \item The statistical methodology (Texas Sharpshooter test with Bonferroni correction) is rigorous and clearly presented, yielding $p < 0.0001$.
    \item The central question---whether the arithmetic structure of perfect numbers constrains optimal computation---is of broad mathematical interest and suitable for the Monthly's expository tradition.
\end{enumerate}

This manuscript has not been published or submitted elsewhere. All code and data are publicly available at \texttt{https://github.com/need-singularity/anima}.

I used Claude (Anthropic) as a writing assistant for formatting and translation. All mathematical content, experimental results, and analysis are my own original work.

Thank you for your consideration.

\closing{Sincerely,}

\end{letter}
\end{document}
